\newpage
\phantomsection
\label{prf:two-element-system-is-field}
\LRAProofFor{thm:two-element-system-is-field}

\begin{remark*}[Return]
\hyperref[thm:two-element-system-is-field]{Return to Theorem}
\end{remark*}

\begin{theorem*}[The Two-Element System Is a Field]
The two-element number system \(\mathbb{F}_2\) is a field.
\end{theorem*}

\begin{proof}[Professional Standard Proof]
\LRAProofBodyStart
\textbf{Professional Standard Proof.}
The operation tables show closure. Addition is associative and commutative by
checking the eight possible triples, \(0\) is the additive identity, and each
element is its own additive inverse. Multiplication is associative and
commutative by the same finite check, \(1\) is the multiplicative identity,
and the only nonzero element is \(1\), whose multiplicative inverse is \(1\).
The distributive law is verified on the eight possible triples. Hence the
field axioms hold.
\end{proof}

\begin{proof}[Detailed Learning Proof]
\LRAProofBodyStart
\textbf{Detailed Learning Proof.}
Because \(\mathbb{B}_2=\{0,1\}\), every universal law involving three
variables can be checked on the eight triples in \(\mathbb{B}_2^3\). The two
tables define outputs only in \(\{0,1\}\), so both operations are closed.
The addition table has \(0+_2x=x=x+_20\) and
\(x+_2x=0\) for \(x=0,1\), giving an additive identity and additive inverses.
The multiplication table has \(1\cdot_2x=x=x\cdot_21\), and the only nonzero
element \(1\) has inverse \(1\). Direct table checking gives associativity,
commutativity, and distributivity, so the structure is a field.
\end{proof}

\begin{remark*}[Proof structure]
Finite operation tables turn the field axioms into a finite verification.
\end{remark*}

\begin{dependencies}
\begin{itemize}
  \item \hyperref[def:two-element-number-system]{Two-Element Number System}
  \item \hyperref[def:two-sided-inverse]{Two-Sided Inverse}
\end{itemize}
\end{dependencies}

\clearpage
